\section{Concurrent Operational Game Semantics}
\label{s:cogs}

In order to establish an exact correspondence between OGS and the encoding into $\pi$-calculus,
as established in Theorem~\ref{t:opcorrCON3_w}, we had to introduce the \opLTS for the 
$\pi$-calculus.

In this section, we explore another way to derive an exact correspondence between OGS and \piI,
by relaxing the Alternating LTS for OGS to allow multiple terms in 
configurations to run concurrently. 
We refer to the resulting OGS as the \emph{Concurrent OGS}, briefly \COGS (we
recall that  \AOGS refers to the Alternating  OGS of Section~\ref{s:AOGS}).  

\subsection{The Concurrent OGS LTS}
\label{ss:cs_cogs-lts}


We introduce \emph{running terms}, ranged over by $A,B$,
as partial maps with finite domains 
from continuation names to $\lambda$-terms.
Each mapping $\pmap{p}{M}$ of a running term can be seen as
a thread evaluating $M$, with $p$ an identifier
used to communicate the result of the evaluation.
%
A \emph{concurrent configuration} is a triple $ \conf{A,\gamma,\phi}$ 
of a running term $A$, an environment 
$\gamma$, and a set of names $\phi$.
Moreover, the domains of $A$ and $\gamma$
must be disjoint.
Both the running term and the environment may be empty.
The continuation structure of such a concurrent configuration
is then defined as the pair $(\dom{A},\cs(\gamma))$
formed by the domain of $A$ and the continuation structure of $\gamma$.
%The available continuation names of a concurrent configuration
%are then defined as the ones in the domain of $A$ plus
%the ones in the continuation structure of $\gamma$.
We extend the definition of the polarity function,
considering names in the domain of both $A$ and $\gamma$ as Player names. 
%

% 
% It is required, and will be implicitely assumed so in the remainder, that, 
% in a configuration, 
% the \Support of
% its running term  and and the flat \Support of its environment are
% disjoint.  

% The \Support of $ \conf{A,\gamma,\phi}$ is the union of the
% \Supports of $A$ and  
% $\gamma$; as usual its flat \Support is the set of names in the \Support.

Passive and active configurations
can be seen as special cases of \COGS configurations with empty and singleton
running term, respectively.
For this reason we still use $\FF, \GG$ to range
over \COGS configurations. Moreover, we freely take  \AOGS
configurations to be  \COGS configurations, 
and conversely for \COGS configurations
with empty and singleton running term, 
omitting the obvious syntactic coercions. 

\begin{figure*}
 \[\begin{array}{l|lll@{}l}
   (P\tau) & \conf{A \cdot \pmap{p}{M},\gamma,\phi} & \arrC{\ \tau \ } & 
     \conf{A \cdot \pmap{p}{N},\gamma,\phi} & \text{ when } M \redv N\\
%
   (PA) & \conf{A \cdot \pmap{p}{V},\gamma,\phi} & \arrC{\ansP{p}{x}} & 
     \conf{A,\gamma \cdot \pmap{x}{V},\phi \uplus \{x\}} \\
%
   (PQ) & \conf{A \cdot \pmap{p}{E[xV]},\gamma,\phi} & \arrC{\questPV{x}{y}{q}} & 
     \multicolumn{2}{l}{\conf{A,\gamma \cdot \pmap{y}{V} \cdot \pmap{q}{(E,p)},
     \phi \uplus \{y,q\}}} \\
%
   (OA) & \conf{A,\gamma\cdot\pmap{p}{(E,q)},\phi} & \arrC{\ansO{p}{x}} & 
     \conf{A \cdot \pmap{q}{E[x]},\gamma,\phi \uplus \{x\}} \\
%
   (OQ) & \conf{A,\gamma,\phi} & \arrC{\questOV{x}{y}{p}} & 
     \conf{A \cdot \pmap{p}{V y},\gamma,\phi \uplus \{y,p\}} 
     & \text{ when } \gamma(x) = V\\
   (IOQ) & \initconf{\phi}{M} & \arrC{\questOinit{p}} &
     \conf{\pmap{p}{M},\emptymap,\phi \uplus \{p\}}
  \end{array}\]
\caption{The LTS for the Concurrent OGS}
\label{fig:conc-ogs-lts}
\end{figure*}
 
We present the rules of \COGS in Figure~\ref{fig:conc-ogs-lts}.
Since there is no more distinction between passive and active configurations,
a given configuration can perform both %different (non-equivalent) 
Player and Opponent actions.
Notice that only Opponent can add a new term to the running term $A$.
% Player does not have a 
% Player has no access to a ``fork'' operation to 
% Notice that we keep garbage-collecting evaluation contexts when performing Opponent Answers.

%Of particular importance in the rest of the paper are
A  \emph{singleton configuration} is a configuration $\FF$
that has only one P-name (that is, in \COGS,  $\FF$ is either of the form 
$\conf{\pmap{p}{M},\emptymap,\phi}$, or  
$\conf{\emptymap,\pmap{x}{V},\phi}$, or $\conf{\emptymap,\pmap{p}{(E,q)},\phi}$).

In this and in the following section $\FF,\GG$ range over \COGS configurations, as reminded
by the index `$\rmm{C}$'
  in the symbols for LTS and behavioural equivalence with  which $\FF,\GG$ appear (e.g.,
  $\genTE C$).


\subsection{Comparison between \COGS and \piI}
\label{ss:compa_cogs_piI}

The encoding of \COGS into \piI is a simple adaptation of that for \AOGS.
We only have to consider the new syntactic element of \COGS, 
namely running
terms; the encoding remains otherwise the same.
% Encoding of environments is as before, namely: 
% \[ 
% \begin{array}{rcl}
% \encoEnvV{\pmap{y}{V}\cdot\gamma  } & \defi & \encoVVa {V} y | \encoEnvV{\gamma  }
% \\[\mypt] 
% \encoEnvV{\pmap{q}{(E,p)}\cdot\gamma  } & \defi & 
% \inp qx .\encoVa {E[x]} p  | \encoEnvV{\gamma  } \\[\mypt]
% \encoEnvV{\emptymap} & \defi & \nil
% \end{array} 
% \]
The encoding of a running term is: 
\[ 
\begin{array}{rcl}
\encoEnvI{\pmap{p}{M}\cdot A  } & \defi & 
 \encoIa {M}{p} |\encoEnvI{ A  } \\
\encoEnvI{\emptymap} & \defi & \nil
\end{array} 
\]
% \[ 
% \begin{array}{rcl}
% \encoEnvV{\pmap{p}{M}\cdot A  } & \defi & 
%  \encoVa {M}{p} | \encoEnvV{ A  } \\
% \encoEnvV{\emptymap} & \defi & \nil
% \end{array} 
% \]
The encoding of configurations is then defined as:
\[
\encoEnvI{\conf{A, \gamma, \phi}}
\defi \encoEnvI A | \encoEnvI \gamma
\]

The results about the operational correspondence between \COGS and \piI 
are as those between \AOGS and \piI under the \opLTS. 
\DSOCTb
%We only report the correspondence for weak transitions. 
We  report the main result, namely the correspondence for weak transitions, as well as 
two lemmas about the  correspondence for strong transitions, since they are also needed
when lifting up-to techniques from 
 \piI to \COGS, in Section~\ref{ss:upto_games}.
The reasoning in their proofs is similar to that 
in the operational  correspondence for \AOGS
(Lemmas~\ref{l:opcorrCON_str} and
\ref{l:opcorrCON2}). 

 
\begin{lemma}[from \COGS to \piI, strong transitions]
\label{l:opcorrCOGS}
$ $ \begin{enumerate}

\item If $\FF \longrightarrowp \FF'$, then 
$ \encoConV \FF \longrightarrowPI  \contr \encoConV {\FF'}$;


\item If $\FF \arrC\ell \FF'$, then
$ \encoConV \FF \arrPI\ell  
\wbPI
\encoConV {\FF'}$.
\end{enumerate} 
\end{lemma} 

\begin{lemma}[from \piI to  \COGS, strong transitions]
\label{l:opcorrCOGS2}
$ $ \begin{enumerate}


\item If 
$ \encoConV \FF \longrightarrowPI \PP$ then there is $\FF'$ such that 
$\FF \longrightarrowp \FF'$  and  
$\PP   \contr \encoConV {\FF'}$ ; 


\item If
$ \encoConV \FF \arrPI\ell \PP$
 then there is $\FF'$ such that
 $\FF \arrC\ell \FF'$
and $\PP \wbPI \encoConV {\FF'}$.

\end{enumerate} 
\end{lemma} 
 

\DSOCTe

\begin{theorem}
\label{t:COGS_pi}
$ $ \begin{enumerate}

\item
If   $\FF \Longrightarrowp \FF'$ then 
$\encoEnvI \FF \LongrightarrowPI  
\contr
\encoEnvI {\FF'}$;

\item
if   $\FF \ArrC{\ell} \FF'$ then 
$\encoEnvI \FF \ArrPI\ell  
\wbPI
\encoEnvI {\FF'}$;

\item %the  converse of (1), i.e.
if 
$\encoEnvI \FF \LongrightarrowPI   \PP$ then there is $\FF'$ such that 
   $\FF  \Longrightarrowp  \FF'$ and
$\PP  \wbPI
\encoEnvI {\FF'}$;

\item %the  converse of (2), i.e.
if 
$\encoEnvI \FF \ArrPI\ell   \PP$ then there is $\FF'$ such that 
   $\FF  \ArrC{\ell} \FF'$ and
$\PP \wbPI
\encoEnvI {\FF'}$.
 \end{enumerate} 
\end{theorem} 

% \finish{make sure $ \ArrC{\widehat{\act}}$ has been introduced (if
% used only here, it can 
% \tr  be done here)} 




\begin{corollary} %[full abstraction for  traces, in the ordinary LTS]
\label{c:ogsp_pi}
For any  \COGS configuration $\FF$ 
 and trace $\tr$,  
we have  $\FF \ArrC{\tr}$ iff $\encoEnvI \FF \ArrPI \tr $.
\end{corollary} 

\noindent
From Corollary~\ref{c:ogsp_pi} 
and Theorem~\ref{t:COGS_pi}, we derive:

\begin{lemma} %[full abstraction for  traces and bisimilarity] %  ordinary LTS]
\label{l:ogsp_pi_traces_bis}
For any $\FF,\FF'$ we have:
\begin{enumerate}
\item
 $\FF_1 \TEp \FF_2$ iff 
 $\encoEnvI {\FF_1} \TEPI \encoEnvI {\FF_2}$;  
\item
 $\FF_1 \wbC \FF_2$ iff 
 $\encoEnvI {\FF_1} \wbPI  \encoEnvI {\FF_2}$.  
\end{enumerate} 
\end{lemma} 

\noindent
To derive the full analogue of Corollary~\ref{c:OGS_opLTSpiI},  
we now show that, on the \piI representation of $\lambda$-terms,
trace equivalence is the same as bisimilarity.
This result needs a little care:
it is known that on deterministic LTSs
bisimilarity coincides with trace equivalence.
\iflong
  (In a deterministic
 LTS, intuitively,  for any process and action there exists at most
 one proof of derivation of that action from the process).
\fi
However,
 the behaviour of the \piI representation of a  \COGS
 configuration
 need not be
deterministic, because there could be
multiple silent transitions as well as multiple output transitions 
(for instance, in \COGS
rule OQ may be
applicable to different terms).

\iflong
\finish{some more explanations could go here. The problem is really
the outputs. with only silent actions a weaker form of determinacy
would make things ok} 
\fi

\begin{lemma}
\label{l:tr_bisi_sing}
For any $M,N$ we have: $
{\encoEnvI M} \TEPI {\encoEnvI N}
$ iff 
${\encoEnvI M} \wbPI {\encoEnvI N}$.
\end{lemma} 

\begin{proof}
Bisimilarity implies trace equivalence. 
 The delicate direction is the converse one.
It is easier to work with the optimised encoding  $\qencV$.
The relation 
\[ 
(\encoConV \FF,  \encoConV \GG) \st \mbox{ $\FF,\GG$ are singleton with $\FF \TEp
\GG$}
\]
is a bisimulation up-to context and up-to $(\contr,
 \wbPI)$. 
The proof uses
Lemmas~\ref{l:immediate_tr} and \ref{l:par_tr}
and is reported in
Appendix~\ref{a:aux_tr_bisi_sing}.
\end{proof}

\noindent
 We can finally  combine Lemmas~\ref{l:tr_bisi_sing} and~\ref{l:ogsp_pi_traces_bis}  
 to  derive that the \COGS and \piI semantics of $\lambda$-calculus
coincide, both for traces and for bisimilarity. 

\begin{corollary}
\label{c:fa_all}
For all $M,N$ we have: 
$\conf{M} \TEp \conf {N}$ iff 
$\conf{M} \wbC \conf {N}$ iff 
$\encoEnvI {{M} } \TEPI\encoEnvI {{N}}$
iff 
$\encoEnvI {{M}} \wbPI \encoEnvI {{N}}$. 
\end{corollary} 

%More details on proofs may be found in Appendix~\ref{a:aux_tr_bisi_sing}.

\subsection{Tensor Product}
\label{ss:tp}
We now introduce a way of combining configurations, which corresponds to the notion of
tensor product of arenas and strategies in (denotational) game semantics.

\begin{definition}
Two concurrent configurations $\FF,\GG$ are said to be \emph{compatible} if 
their polarity functions  $\Label_{\FF},\Label_{\GG}$ are compatible~--- that is,
for all $a \in \dom{\Label_{\FF}} \cap \dom{\Label_{\GG}}$,
we have $\Label_{\FF}(a) = \Label_{\GG}(a)$.
\end{definition}

\noindent
Following Definition~\ref{d:interact},
two $\pi$-calculus processes $P_1,P_2$ \emph{cannot interact} 
if there is no name that appears free in both processes
and with opposite polarities 
(that is, in input position in one, 
and output position in the other). 
As
 in \piI all names exchanged are fresh, 
the  `cannot interact' property is preserved by transitions.

\begin{lemma}
  \label{l:ci_trans}
If $P_1,P_2$ cannot interact and 
$P_1 | P_2 \ArrPI \trace P $, then $P =  P_1'| P_2'$, for some 
$P_1', P_2'$ that cannot interact.
\end{lemma}
 
\noindent
The \piI translation of
two compatible \COGS configurations gives two processes that 
cannot interact with each other.

\begin{lemma}
  \label{l:compat_cannot-interact}
  Taking two \COGS configurations $\FF_1,\FF_2$,
  if their polarity functions are compatible then $\encoConI {\FF_1}$ and
  $\encoConI {\FF_2}$ cannot interact.
\end{lemma} 


\begin{definition}
  \label{d:tensor_cogs}
For  compatible configurations
$ \FF = \conf{ A , \gamma, \phi }$
and 
$ \GG = \conf{ B , \delta, \phi' }$,  the \emph{tensor product} 
$\FF \appendSYMB \GG$ is  defined as 
\[ \appendCon \FF \GG  \defeq
 \conf{ A \cdot B , \gamma \cdot \delta, \phi \cup \phi' }
 \] 
 \end{definition}

\noindent
 The polarity function of $\appendCon \FF \GG$ 
 is then equal to $\Label_{\FF} \cup \Label_{\GG}$, and
$\encoConI {
\appendCon
{\FF_1} {\FF_2}}
\equiv 
\encoConI {
\FF_1}
|
\encoConI {
 {\FF_2}}$,  where $\equiv$ is the standard structural congruence of  $\pi$-calculi.
 % \ds{above:    defined?}
 % \gj{I guess using = is not correct because of the need of 
 % associativity and commutativity reasoning over parallel composition ?}
 

In the following, we write $\interl{\trace_1}{\trace_2}$ for the set of traces
obtained from an interleaving of the elements in the sequences $\trace_1$ and $\trace_2$.

We relate the traces generated by the tensor product 
of two configurations to the set of interleavings of traces generated 
by the component configurations themselves. 
The result is proved by going through the $\pi$-calculus,
via the following lemma.


\begin{lemma}
  \label{l:dec_pi}
  Suppose $P_1,P_2$ cannot interact.
  \begin{enumerate}
  \item
  If $P_1 | P_2 \ArrPI \trace   P_1'| P_2'$,
  then, for $i=1,2$,  there is  $ \trace_i $ 
  such that $P_i \ArrPI { \trace_i} P_i'$, and $ \trace  
  \in   \interl { \trace_1   }{ \trace_2 }$;
  \item 
  Conversely,
  if for  $i=1,2$,  we have $P_i \ArrPI { \trace_i} P_i'$, and $ \trace  
  \in   \interl { \trace_1   }{ \trace_2 }$, then 
  $P_1 | P_2 \ArrPI \trace   P_1'| P_2'$.
  \end{enumerate} 
\end{lemma} 


\begin{proof}
  The interesting case is (1). 
  Consider the sequence of one-step  actions (including silent actions) that 
  are performed to obtain the trace
  $P_1 | P_2 \ArrPI \trace   P_1'| P_2'$.  Then tag each of these actions with ${\tt L}$ or 
  ${\tt R}$ depending on whether in the derivation proof of that action, the last rule
  applied was {\tt parL} or  {\tt parR}.
  Then $ \trace_1$ is obtained by collecting the subsequence of $\trace$ with  
  the {\tt ParL} tag, and similarly for   $ \trace_2$ with
  {\tt ParR}.
\end{proof}

\noindent
We finally prove that the set of traces generated by $\FF_1 \otimes \FF_2$
is the  union of the sets of interleaving $\interl{\trace_1}{\trace_2}$,
for $\FF_1 \ArrC{\trace_1}$ and $\FF_2 \ArrC{\trace_2}$.
\begin{lemma}
\label{l:COGS-interleave-tensor}
%Suppose $\FF_1,\FF_2$ are compatible concurrent configurations.
%The set of traces generated by $\FF_1 \otimes \FF_2$
% is the  union of the sets of interleaving $\interl{\trace_1}{\trace_2}$,
%for $\FF_1 \ArrC{\trace_1}$ and $\FF_2 \ArrC{\trace_2}$.
\iflong
Suppose $\FF_1,\FF_2$ are compatible concurrent configurations.
\begin{enumerate}
\item\label{l:COGS-interleave-tensor_1}
If $\appendCon
 {\FF_1} {\FF_2} \ArrC \trace  \FF  $, then, 
for $i=1,2$,  there are traces $ \trace_i $
such that $ {\FF_i} \ArrC { \trace_i} \FF'_i $, and $ \trace  
\in   \interl { \trace_1   }{ \trace_2 }$, and 
$\FF = \appendCon
 {\FF'_1} {\FF'_2}$;
\item\label{l:COGS-interleave-tensor_2}
Conversely,
if for  $i=1,2$, we have
 $ {\FF_i} \ArrC { \trace_i} \FF'_i$,
and
$ \trace  
\in   \interl { \trace_1   }{ \trace_2 }$, 
then 
$ 
\appendCon
{\FF_1} {\FF_2} \ArrC \trace   \appendCon
 {\FF'_1} {\FF'_2}$.
\end{enumerate}
\fi
\end{lemma}  

\begin{proof}
  It follows from Lemma~\ref{l:dec_pi}, 
  Lemma~\ref{l:compat_cannot-interact},
  Corollary~\ref{c:ogsp_pi}, and the fact that
  $\encoConI {
  \appendCon
  {\FF_1} {\FF_2}}
  \equiv 
  \encoConI {
  \FF_1}
  |
  \encoConI {
  {\FF_2}}$.
\end{proof}

\noindent
The tensor product of \AOGS configurations will be introduced
in Section~\ref{a:tensor-aogs}.
%is defined similarly,
%with the additional hypothesis that at most one of the two configurations is  active,
%in order for their tensor product to be a valid \AOGS configuration.
%The details can be found in Appendix~\ref.



\subsection{Up-to techniques for games}
\label{ss:upto_games}

We introduce up-to techniques for \COGS, which allow, in bisimulation proofs, 
to split two \COGS configurations  
into separate components and then to reason separately on these.
These up-to techniques are directly imported from \piI.
Abstract settings for up-to techniques have been developed, see
\cite{SanPous,PousS19}; we cannot however derive the OGS techniques from
them because these settings are specific to first-order LTS (i.e., CCS-like, 
without binders within actions). 
%
They are then used to prove that
\COGS and \AOGS yield the same equivalences on
$\lambda$-terms.
A further application is in Section~\ref{s:enf}, 
discussing eager normal-form bisimilarity. 

In order to compare two configurations,
they must share the same structure, as defined by the following notion.
\begin{definition}
\label{d:compa}
Two configurations $\FF,\GG$ are \emph{support-equivalent}, written 
$\okC \FF\GG$, when they have the same polarity function (and so the same names) and the
same continuation structure.

\noindent
A relation $\R$ on configurations is \emph{well-formed} if it relates
support-equivalent configurations. 
\end{definition} 

\noindent
Below, all relations on configurations are meant to be well-formed. 
 %
Given a well-formed relation $\R$ we write
\begin{itemize}
\item $\appendRel \R $ for the relation
$
\{ (\FF_1,\FF_2) \st
\mbox{$\exists$ $\GG$ s.t.\ }
 \FF_i = \appendCon{\FF'_i}\GG \mbox{ ($i=1,2$)}  \mbox{ and } 
 \FF'_1\RR \FF'_2\}
$.
\item
$\uptoComp \R $ for the reflexive and transitive closure of 
$\appendRel \R $.

Thus from $\FF_1 \RR  
% \uptoComp{\R} 
\GG_1$ and  $\FF_2 \RR
% \uptoComp{\R} 
\GG_2$ we obtain 
$(\appendCon{\FF_1}{\FF_2})  \uptoComp{\R}   
(\appendCon{\GG_1}{\FF_2}) \uptoComp{\R}  \appendCon{\GG_1}{\GG_2}.$
% This decomposition will be used to lift relations over plain or singleton configurations,
% into relations over concurrent configurations.



\item 
$\uptoCompRed \R $ for the closure of $\uptoComp \R $ under
reductions. That is, $\FF_1 \uptoCompRed \R \FF_2$ holds if there are 
$\FF_i'$, $i=1,2$ with $\FF_i \Longrightarrowp \FF'_i$ and 
$\FF_1' \uptoComp \R \FF_2'$.
(As $\Longrightarrowp$ is reflexive, we may have $\FF_i =
\FF'_i$.) 
\end{itemize}
 
\begin{definition}
\label{d:uptoComp}
A relation $\R$ on configurations is a \emph{bisimulation up-to
 reduction and   composition } if whenever $\FF_1 \RR \FF_2$: 
\begin{enumerate}
\item
 if $\FF_1 \arrC\act \FF_1'  $ then there is $\FF_2' $ such that 
$\FF_2 \Arcapp\act \FF_2'  $ and 
$\FF'_1 \uptoCompRed \R \FF'_2$ ;

\item the converse, on the transitions from  $\FF_2$.
\end{enumerate}
\end{definition}


\DSOCTb

\noindent
Sometimes we do not need  the `closure under reduction'; that is,  referring to  the clauses
of Definition~\ref{d:uptoComp} 
 we may take  
$\FF'_i = \FF_i$, ($i=1,2$).
In this case  
we say that $\R$ is a 
`{bisimulation up-to
   composition}'.
%  and just take, in clause (1) above   (and similarly for
% its converse (2)) 
% $\FF'_1 \uptoComp \R \FF'_2$.

For the soundness of 
bisimulation up-to reduction and composition,
we first need the following Lemma~\ref{l:uptoOGSpi}, 
which is a  consequence of the compositionality of the mapping from
\COGS to \piI.  

\begin{lemma}
\label{l:uptoOGSpi}
For a relation $\R$ on \COGS configurations, 
if 
${(\FF_1, \FF_2)} \in  { \uptoComp \R}$, then 
we have $( \encoConV {\FF_1},  \encoConV {\FF_2}) \in  \ctx{( \encoConV{ \R })}$. 
\end{lemma} 

\noindent
We are now ready to show  the soundness of 
bisimulation up-to reduction and composition.

\begin{theorem}
\label{t:uptoComp}
If $\R$ is bisimulation up-to
reduction and composition then ${\R} \subseteq {\wbC}$.
\end{theorem} 


\begin{proof}
We prove the result via \piI, exploiting Lemma~\ref{l:ogsp_pi_traces_bis}(2).
Here again, it is convenient to use
the optimised encoding $\qencV$.

If $\R$ is a relation on \COGS configurations, then 
$\encoConV \R$ is the relation on \piI processes obtained by mapping each pair
in $\R$ in the expected manner: 
\[ \encoConV \R \defi \{ (\encoConV {\FF_1}, \encoConV {\FF_2}) \st \FF_1 \RR \FF_2 
 \} \] 
And, if $\S$ is a relation on \piI processes, then $\ctx \S$ is the closure of $\S$ under
polyadic contexts, i.e., the set of pairs  of the form  
$(\ct {P_1,\ldots, P_n},\ct {Q_1,\ldots, Q_n})$ where $\qct$ is a polyadic context and for
each $i$, $P_i \SS Q_i$. 

As a consequence of the compositionality of the mapping from
\COGS to \piI,
for  a relation   $\R$ on \COGS configurations, 
if ${(\FF_1, \FF_2)} \in  { \uptoComp \R}$, then 
$( \encoConV {\FF_1},  \encoConV {\FF_2}) \in  \ctx{( \encoConV{ \R })}$. 




Now, let $\R$ be the relation in the hypothesis of the theorem.
We  prove the theorem by  showing that 
$\encoConV \R$ is a bisimulation
 up-to context and up-to $(\contr,
 \wbPI)$ in 
\piI, and then appealing to Theorem~\ref{t:up-to-con-expa-nointer}.  

Suppose 
$\encoConV {\FF_1}  \arrPI \ell P$.  Then
there is $\FF_1'$ with  
$\FF_1  \arrPI \ell \FF_1' $ and $P  \wbPI \encoConV {\FF'_1}$ (this is derived from
Lemma~\ref{l:opcorrCOGS2}).  
Since $\R$ is bisimulation up-to
 reduction and  composition and $\FF_1 \RR \FF_2$, there are 
  $\FF_1''$, and $\FF_2'$, with  
$\FF'_1 \Longrightarrowp \FF''_1$, 
$\FF_2
  \arrPI \ell \Longrightarrowp \FF_2' $ and
$\FF''_1 \uptoComp \R \FF'_2$. 


Using  Theorem~\ref{t:COGS_pi}(1) 
%Lemma~\ref{l:Redp} 
from 
$\FF'_1 \Longrightarrowp \FF''_1$ and the inclusion ${\contr} \subseteq
{\wbPI}$, 
we derive
$\encoConV {\FF'_1}  \wbPI  \encoConV {\FF''_1}$; 
and  by   
Theorem~\ref{t:COGS_pi}(1),
%Lemmas~\ref{l:Redp}, 
and  Lemmas \ref{l:opt_sound} and \ref{l:opcorrCOGS},  from 
$\FF_2
  \arrPI \ell \Longrightarrowp \FF_2' $ 
we derive
$\encoConV {\FF_2}  \arrPI \ell \wbPI  \encoConV {\FF'_2}$.   


Finally,  
\iflong
by Lemma~\ref{l:uptoOGSpi},  
\fi
from $(\FF''_1, \FF'_2) \in  \uptoComp \R$, we derive 
 $( \encoConV {\FF'_1},  \encoConV {\FF'_2}) \in  \ctx{( \encoConV{ \R })}$. 
This closes the proof, up-to polyadic contexts and bisimilarity.

The case when 
$\encoConV {\FF_1}$ makes a silent move is simpler,  
using expansion in place of bisimilarity.
\end{proof} 

\noindent
We now present an additional up-to technique, 
which will be used when relating 
OGS bisimilarity with eager normal form bisimilarity in Section~\ref{s:enf}.
When $\R$ is a relation on singleton configurations, 
a straightforward simplification 
of `bisimulation up-to reduction and composition'
allows us to play the bisimulation game only on visible actions;
the bisimulation clause becomes: 
\begin{itemize}
\item
 if $\FF_1 \Longrightarrowp\arrC\ell \FF_1'  $ then there is $\FF_2' $ such that 
$\FF_2\Longrightarrowp\arrC\ell \FF_2'  $ and 
$\FF'_1 \uptoComp \R \FF'_2$.
\end{itemize} 
We call this a \emph{singleton bisimulation up-to composition}; 
its soundness is a corollary of that for bisimulation   
up-to reduction and  composition, 
derived from the up-to techniques for \piI and the
mapping from \COGS  onto \piI. 



\begin{proposition}
\label{p:sing_bis}
If $\R$ is a singleton bisimulation up-to composition, then ${\R} \subseteq {\wbC}$.
\end{proposition} 

\begin{proof}
 Take 
$$\S \defi \R \cup \{ (\FF_1,\FF_2) \st 
%\begin{array}[t]{l}
\mbox{ there are } \FF'_1, \FF'_2  \mbox{ with } \FF'_1 \RR \FF'_2 
\mbox{ and   }\FF'_i \Longrightarrowp \FF_i  \mbox{ $i=1,2$}  \}
%\end{array} 
$$
Intuitively,  $\S$  is the closure of $\R$ under deterministic reductions. 
Then $\S$ is a bisimulation up-to composition.
Note that we exploit the fact that transitions of singleton configurations are deterministic.
\end{proof} 


\DSOCTe

\begin{remark}
\label{r:CONFnoSINGL} 
Results such as Corollary~\ref{c:fa_all} and 
Theorem~\ref{t:uptoComp} might suggest that the equivalence between
two configurations implies the equivalence of all their singleton 
components. That is, if $\FF \TEp \GG$, with 
$\pmap{p}{M}$ part of $\FF$ and $\pmap{p}{N}$ part of $\GG$, 
then also $\conf {\pmap{p}{M}} \TEp \conf {\pmap{p}{N} }$. 
A counterexample is given by the configurations 
$$
\begin{array}{rcl|l}
\FF_1 & \defi & \conf{ \pmap{p_1}{M}\cdot \pmap{p_2}{\Omega} } 
& \text{ where } M \defi (\lambda z. \Omega)(x \lambda y . \Omega) \\
\FF_2 & \defi & \conf{ \pmap{p_1}{\Omega}\cdot \pmap{p_2}{M} }
& 
\end{array}
 $$ 
Intuitively the reason why $\FF_1 \TEp \FF_2$  is that
the term $M$ can produce an output (along the variable $x$), 
 but an observer  will
never obtain access to  the  name at which $M$ is located
($p_1$ or $p_2$).
That is, the term $M$ can interrogate $x$,
but it will never answer,  neither at $p_1$ nor at $p_2$.
%
However the result does hold when $\FF,\GG$ are reachable from
$\lambda$-terms, that is, there are $M,N$ and a trace $s $ such
that $\conf M \ArrC s  \FF$ and $\conf N \ArrC s \GG$. 
\end{remark} 

\begin{remark}
  \label{rm:up-to-comp-gs}
  Up-to composition can be understood as a decomposition principle for configurations 
  whose components do not interact with each other, even though they may share names 
  that are only used for communication with the environment. 
  This decomposition reduces the OGS bisimulation game to the analysis of singleton configurations, 
  effectively allowing one to reason independently on each branch of a (possibly infinite) 
  Böhm-like trees representation of normal forms, as it will be
  made apparent in Section~\ref{s:enf} with Lassen trees. 
  Through the well-known correspondence between such branches and P-views in game semantics~\cite{curien1998abstract}, 
  this yields a view-based reasoning principle: up-to composition enables one to analyse strategies 
  by considering each P-view independently, in a forward manner. 
  In this sense, it captures and operationalises a characteristic feature of innocent strategies, 
  namely that behaviours can be decomposed along views without mutual interference between them.
\end{remark}


\subsection{Relationship between Concurrent and Alternating OGS}
\label{subsec:relate-conc-alt-lts}


In Section~\ref{s:encoGames} we have proved that the
trace-based and bisimulation-based semantics produced by
\AOGS and by \piI under the \opLTS coincide. 
In Section~\ref{ss:compa_cogs_piI} we have obtained the
same result for \COGS and \piI under the ordinary LTS.
In this section, we develop these results to conclude that
all such equivalences for $\lambda$-terms actually coincide.  In other words, the equivalence
induced on $\lambda$-terms by their representations in OGS and \piI is
the same, regardless of whether we adopt the alternating or concurrent
flavour for OGS, the \opLTS or the ordinary LTS in \piI, a trace or a
bisimulation semantics.

For this, in one direction we show that the trace semantics induced by
\AOGS coincides with the subset of \emph{alternating} traces of \COGS.
 
\begin{lemma}
\label{l:traces_AOGS_COGS}
For any \AOGS configuration $\FF$, 
$\FF \ArrA{\tr}$ if and only if $\FF \ArrC{\tr}$ with $\tr$ an alternating
trace starting with an output when $\FF$ is active.
\end{lemma} 

\begin{proof}
Taking a trace $\tr$ such that $\FF \ArrA{\tr}$,
we directly get that $\tr$ is alternating.
Moreover, $\tr$ starts with an output when $\FF$ is active.
Since the rules of \AOGS LTS are a subset of the ones of \COGS,
we also get that $\FF \ArrC{\tr}$.

In the opposite direction, 
we take $\tr$ an alternating trace such that 
$\FF \ArrC{\tr}$ and $\tr$ starts 
with an output when $\FF$ is active.
We prove $\FF \ArrA{\tr}$ by induction over $\tr$.
If the trace is empty, this is direct.
Otherwise, we write $\tr$ as $\act \tr'$ such that
$\FF \arrC{\act} \GG \ArrC{\tr'}$.
Then we reason by case analysis over $\FF$:
\begin{itemize}
  \item either $\FF$ is active so that $\act$ must be an output by the statement
  of the lemma.
  So we get that $\FF \arrA{\act} \GG$ and $\GG$ is passive.
  \item or $\FF$ is passive, so that $\act$ has to be an input.
  We get that  $\FF \arrA{\act} \GG$ with $\GG$ active.
  Since $\tr$ is alternating, we then get that $\tr'$ is
  starting with an output. 
\end{itemize}
Since $\tr'$ is alternating,
in both cases we conclude using the induction hypothesis.
\end{proof}
 
\noindent
In the opposite direction, the
difficulty is that the \COGS LTS can perform more actions than the \AOGS:
it can perform input actions on active configurations.
So we state the correspondence only on singleton configurations.
But neither the \COGS nor the \AOGS output transitions preserve the fact of
being singleton.
The solution is to use bisimulation up-to composition over the \COGS LTS
in order to keep reasoning on singleton configurations.
We then have to only extract singleton configurations
from \AOGS configurations that can effectively be triggered directly by an action.

\begin{definition}
A singleton \COGS configuration $\FF$ is 
a \emph{triggerable} subconfiguration of an \AOGS configuration $\FF'$ when:
\begin{itemize}
  \item either $\FF'$ is active of the shape
  $\conf{M,p,\gamma,\phi}$,
  and $\FF = \conf{\pmap{p}{M},\emptymap,\phi \backslash{\dom{\gamma}}}$
  \item or $\FF'$ is passive of the shape
  $\conf{\gamma,\phi}$,
  and $\FF$ is:
  \begin{itemize}
    \item either of the shape
    $\conf{\pmap{x}{V},\phi \backslash{\dom{\gamma}} \cup\{x\}}$
    with $\pmap{x}{V} \subseteq \gamma$;
    \item or of the shape $\conf{\pmap{p}{(E,q)},\phi \backslash{\dom{\gamma}} \cup\{p\}}$
    with $\pmap{p}{(E,q)} \subseteq \gamma$.
  \end{itemize}
\end{itemize}
\end{definition}
\noindent
Notice that a triggerable subconfiguration $\FF$ of $\FF'$ is not formally a subconfiguration of $\FF'$,
as one needs to remove the Player names of $\FF'$ not in the domain of $\FF$, and in the case of $\FF'$ active,
we transform the singleton running term of $\FF$ into a term and a continuation name.

Configurations $\conf{\pmap{x}{V},\phi}$ and
$\conf{\pmap{p}{(E,q)},\phi}$ are not triggerable 
when $\FF$ is active,
since they cannot be directly triggered by the \AOGS LTS
(while they can be in the \COGS LTS).


\reviewComment{"triggerable subconfiguration" $\rightarrow$ terminology is confusing since
a triggerable subconfiguration is not a subconfiguration?}
\reviewCommentDS{I would add here this sentence:   
 $\FF$ being  
a triggerable subconfiguration of $\FF'$  means that $\FF$ is obtained from $\FF'$ as in
the definition; it 
does not imply that  $\FF$ is a syntactic subterm
of $\FF'$.}
\reviewCommentGJ{I've added an explanation above.}


\begin{lemma}
\label{l:OGS_COGS}
If $\FF, \GG$ are \AOGS singleton configurations and 
$\FF \wbA \GG$, then also $\FF \wbC \GG$. 
\end{lemma}

\begin{proof}
  Let $\R$ be the relation on \COGS singleton configurations with $\FF \RR \GG$ if
   $\okC \FF\GG$ and there exist \AOGS configurations $\FF', \GG'$
    with $\FF' \wbA \GG'$ such that 
    $\FF$ and $\GG$ are triggerable singleton configurations
    respectively of $\FF'$ and $\GG'$.

  First notice that if $\FF, \GG$ are \AOGS singleton configurations and 
  $\FF \wbA \GG$, we directly get that 
  $\FF \RR \GG$ since $\FF$ and $\GG$ are 
  triggerable subconfigurations of themselves,
  once the implicit coercion of \AOGS configurations is applied
  to \COGS configurations.
  
  We now show that $\R$ is a bisimulation up-to composition.
  We take $\FF \RR \GG$, such that there exist $\FF', \GG'$
  as stated by the definition of $\R$.
  Suppose 
  \begin{equation}
  \label{e:KoaH}
   \FF \arrC{\act} \FF_1 
  \end{equation} 
  in \COGS.  
  First, we suppose that $\FF,\GG$ are both passive, 
  so that the rule can be OA or OQ.  
  We assume it is OQ; the case for OA is similar. 
  We thus have, for some $p,V,x,y$ 
  \[\begin{array}{rcl}
   \FF &=&  \conf{ \gamma_1, \phi} \mbox{ with }  \gamma_1 = \pmap{x}{V} \\
   \act & = & \questO{x}{y,p} \\
   \FF_1 & = & \conf{\pmap{p}{V y}, \gamma_1,\phi \uplus \{y,p\}}
  \end{array}\]
  Moreover, we have $\GG = \conf{ \delta_1, \phi}$ 
  with $\delta_1 = \pmap{x}{W}$.
  Since $\FF,\GG$ are triggerable respectively from $\FF',\GG'$,
  we get that both $\FF',\GG'$ are passive.
  From $\FF' \wbA \GG'$, we get that $\okC{\FF'}{\GG'}$,
  so there exist $\gamma, \delta$ sharing the same domain such that:  
  \[\begin{array}{rcl}
  \FF' & = & \conf{ \gamma_1 \cdot \gamma,\phi \uplus \dom{\gamma}} \\ 
  \GG' & = & \conf{ \delta_1 \cdot \delta,\phi \uplus \dom{\delta}} 
  \end{array}\]
  We also have 
  $ \FF' \arrA{\act} \conf{V y,p, \gamma_1 \cdot \gamma, \phi\uplus \dom{\gamma} \uplus \{y,p\}} 
  \defi \FF''$. 
  Since $\FF' \wbA \GG'$, 
  \[ \GG'  \arrA{\act} \conf{W y,q, \delta_1 \cdot \delta, \phi \uplus \dom{\delta} \uplus \{y,p\}}  
  \defi \GG''\]
  and $\FF'' \wbA \GG''$.
  %
  Let $\GG_1 =   \conf{\pmap{q}{W y}, \delta_1,\phi}$.
  We also have 
  \[ \GG \arrC  \act \GG_1. \] 
  This is sufficient to match \reff{e:KoaH}, as 
  $\FF_1 \uptoComp \R  \GG_1$. Indeed we have:
  \begin{itemize}  
  \item $\conf{\gamma_1,\phi\uplus \{y,p\}} \R  \conf{\delta_1,\phi\uplus \{y,p\}}$
  since $\FF \RR \GG$;
  \gj{There is a small glitch: we have extra names $y,p$ here
  that are not in $\FF,\GG$.}
  \item  $\conf{\pmap{p}{V y}, \phi\uplus \{y,p\}} \R  \conf{\pmap{q}{W y}, \phi\uplus \{y,p\}}$
  since they are triggerable subconfigurations respectively of $\FF''$ and $\GG''$.
  \end{itemize}

  \noindent
  Now suppose that $\FF,\GG$ are active,
  so there are 3 possibilities of transitions, corresponding to rules 
  P$\tau$, PA, PQ. The case of 
  P$\tau$ is straightforward since bisimilarity is preserved by internal
  actions.  We only look at PQ, as PA is simpler. 
  Thus we have 
  $\FF = \conf{\pmap{p}{E[x V]},\emptymap,\phi}$, for some $E,x,V$ and: 
  \begin{equation}
  \label{e:KpqH}
   \FF  \arrC{\questPV{x}{y}{q}} 
        \conf{\pmap{y}{V} \cdot \pmap{q}{(E,p)}, \phi \uplus \{y,q\}} 
        \defi \FF_1
  \end{equation} 
  Since $\GG$ shares the same continuation structure as $\FF$,
  we have $\GG = \conf{\pmap{p}{M},\emptymap,\phi}$ for some term $M$.
  Since $\FF,\GG$ are triggerable respectively from $\FF',\GG'$,
  there exist $\gamma,\delta$ sharing the same domain such that:  
  \[\begin{array}{rcl}
    \FF' & = & \conf{E[x V], p, \gamma, \phi\uplus\dom{\gamma}} \\ 
    \GG' & = & \conf{M, p, \delta, \phi\uplus\dom{\delta}} 
  \end{array}\]
  We also have 
  \[
  \FF' \arrA{\questPV{x}{y}{q}} 
   \conf{ \pmap{y}{V} \cdot \pmap{q}{(E,p)} \cdot \gamma, \phi  \uplus \{y,q\}}
   \defi \FF''
  \]
  As $\FF' \wbA \GG'$, there exist $\Ep,W$ such that 
  \[
  \GG'  \Longrightarrow 
   \conf{\Ep[x W],p,\delta,\phi}
  \arrA{\questPV{x}{y}{q}} 
  \conf{\pmap{y}{W} \cdot \pmap{q}{(\Ep,p)} \cdot \delta, \phi \uplus \{y,q\}} 
  \defi \GG''
  \]
  with
  $\FF'' \wbA \GG''$. 
  Defining $\GG_1$ as
  $\conf{\pmap{y}{W} \cdot \pmap{q}{(\Ep,p)}, \phi \uplus \{y,q\}}$,
  we also have $\GG \ArrC{\questPV{x}{y}{q}} \GG_1$. 
  We deduce that both 
  $\conf{ \pmap{y}{V},\phi \uplus \{y\}} 
   \RR
   \conf{ \pmap{y}{W},\phi \uplus\{y\}}$ 
  and 
  $\conf{ \pmap{q}{(E,p)},\phi \uplus \{q\}}
   \RR
   \conf{ \pmap{q}{(\Ep,p)},\phi \uplus \{q\}}$ hold
   since all these singleton configurations
   are triggerable from either $\FF''$ or $\GG''$,
   and that $\FF'' \wbA \GG''$.
  Hence $\FF_1 \uptoComp \R  \GG_1$.
  
  In summary, as an answer to the challenge 
  \reff{e:KpqH}, we have found  $\GG_1$ such that 
   $\GG \ArrC{\questPV{x}{y}{q}}
   \GG_1$ and  
  $\FF_1 \uptoComp \R  \GG_1$; this closes the proof.
\end{proof} 

\noindent
From Corollaries~\ref{c:OGS_opLTSpiI}, \ref{c:fa_all}, and 
Lemmas~\ref{l:traces_AOGS_COGS}, \ref{l:OGS_COGS} we can thus conclude. 
\begin{corollary}
\label{c:OGS_COGS}
For any $\lambda$-terms $M,N$, the following statements are the same: 
$\conf{M} \TEA \conf {N}$;
$\conf{M} \wbA \conf {N}$;
$\conf{M} \TEp \conf {N}$;
$\conf{M} \wbC \conf {N}$;
$\encoEnvI {{M} } \TEN\encoEnvI {{N}}$; 
$\encoEnvI {{M}} \wbOPI \encoEnvI {{N}}$; 
$\encoEnvI {{M} } \TEPI\encoEnvI {{N}}$; 
$\encoEnvI {{M}} \wbPI \encoEnvI {{N}}$. 
\end{corollary} 